\documentclass[11pt,a4paper]{article}

\usepackage[margin=1in]{geometry}
\usepackage{booktabs}
\usepackage{tabularx}
\usepackage{enumitem}
\usepackage{hyperref}

\title{Experimentation Roadmap\\
\large Haptic Ground-Truth Pipeline}
\author{Project Notes}
\date{\today}

\begin{document}
\maketitle

\begin{abstract}
\noindent
Short record of approaches we tried for video-to-haptic ground-truth generation, their outcomes, and whether we changed approach or continued.
\end{abstract}

\section{Goal}
From a short video, make four vibration sound files (algorithms A-D) so people can feel and rate them.
We want continuous flow (like thunder rumble) plus clear hits (like blasts), similar to a dense hand made reference map not only sparse spikes.

\section{Current stack}
\begin{enumerate}[leftmargin=*,itemsep=2pt]
  \item Pull mono sound from the video (44.1\,kHz).
  \item Find candidate moments, then classify them with a sound model (AST) and a video model (ViViT), and merge into an event list (\texttt{events.json}).
  \item Keep only chosen event types (weather, gunshot, explosion, vehicle) and build a filtered audio clip (\texttt{gated\_audio.wav}).
  \item Run Sound2Hap A-D: quiet continuous vibration underneath (bed) plus stronger short hits on top (accents). The four algorithms themselves are unchanged.
\end{enumerate}

\section{Terms}
\begin{itemize}[leftmargin=*,itemsep=3pt]
  \item Fixed list of event types we care about (taxonomy).
  \item Model that links images, audio, and text in one shared space (ImageBind).
  \item Finds objects in each video frame with boxes (YOLO).
  \item Break one long blob of sound into separate bangs at each loud peak (peak split).
  \item Snap the event start to the exact moment the bang begins (onset refine).
  \item Keep audio only at the times of selected event types (temporal + category gate).
  \item Only pass short sharp bangs, ignore long rumble (impulsive gate).
  \item How many real events we find versus miss (detection recall).
  \item Scan every few seconds for steady sounds, and snap bangs using a sudden change in the spectrum (hop + flux onset).
  \item Automatically boost quiet parts so they are not ignored (AGC level).
  \item Quiet continuous vibration that runs under the whole clip (bed).
  \item Mute the quietest share of the track so real silence stays silent (percentile silence).
  \item Squeeze the gap between loud and quiet so soft rumble still feels (envelope-exponent compression).
  \item Average nearby gain values so level changes are soft, not jumpy (mean-smooth compression).
  \item Remember recent loud levels so a bang does not over boost the quiet parts next to it (peak hold).
  \item Scale the quiet continuous layer to a target strength (normalize bed).
  \item Strong short hits placed on top of the quiet layer (accents).
  \item Phone code that turns WAV loudness into vibration strength (parser).
\end{itemize}

\section{Experiments}

\subsection*{Finding what happens in the video}
\begin{center}
\small
\begin{tabularx}{\textwidth}{@{}>{\raggedright\arraybackslash}p{4cm}>{\raggedright\arraybackslash}p{4cm}>{\raggedright\arraybackslash}X@{}}
\toprule
\textbf{What we tried} & \textbf{Outcome} & \textbf{Changed to / Continued?} \\
\midrule
Early notebook with shared image-audio-text model (ImageBind) / video model (VideoMAE) & Incomplete, because it was not linked to Sound2Hap & \textbf{Changed to} frozen sound model (AST) + video model (ViViT) package \\
\addlinespace
Motion-to-vibration idea (HapticLens style) & No named event types (weather, gunshot, etc.) & \textbf{Changed to} fixed event list (taxonomy) + Sound2Hap A-D \\
\addlinespace
Ask ImageBind with text queries for events & Too heavy for Colab flow & \textbf{Changed to} AST + ViViT only \\
\addlinespace
Find objects in frames with boxes (YOLO) for events & Misses off-screen or sound-only cues & \textbf{Continued} AST + ViViT \\
\addlinespace
Colab Google Drive mount & Permission friction for users & \textbf{Changed to} simple file upload \\
\addlinespace
Require both sound and video to accept a gunshot & Missed many explosions & \textbf{Changed to} allow sharp bangs from sound alone, and continued detection \\
\addlinespace
Merge nearby blasts into one long event & Many blasts became one long event & \textbf{Changed to} break at loud peaks (peak split) + snap bang start (onset refine), and continued \\
\addlinespace
Sound2Hap signal-processing A-D & Working candidate generators & \textbf{Continued} (kept as-is, not redesigned) \\
\bottomrule
\end{tabularx}
\end{center}

\subsection*{Choosing which sounds become vibration}
\begin{center}
\small
\begin{tabularx}{\textwidth}{@{}>{\raggedright\arraybackslash}p{4cm}>{\raggedright\arraybackslash}p{4cm}>{\raggedright\arraybackslash}X@{}}
\toprule
\textbf{What we tried} & \textbf{Outcome} & \textbf{Changed to / Continued?} \\
\midrule
Heavy AI noise removal before anything else & Would damage gunshot/explosion edges & \textbf{Changed to} keep audio only at times of selected event types (temporal + category gate) \\
\addlinespace
Pass vehicle rumble and explosions together & Filtered audio dominated by rumble & \textbf{Changed to} only short sharp bangs (impulsive gate) briefly \\
\addlinespace
Only short sharp bangs (impulsive gate) & Missed sustained thunder / tank rumble & \textbf{Changed to} allow weather/vehicle again via taxonomy defaults \\
\addlinespace
One folder of 4 WAVs per detected event & Not desired for evaluation UX & \textbf{Changed to} one full-length WAV per algorithm (4 total) \\
\addlinespace
Skip A-D when no selected events & Empty outputs, Colab shows ``tracks missing'' & \textbf{Continued} by design, because messaging is still unclear \\
\addlinespace
Find loud moments first, then classify only those (propose-then-classify) & Faster, but still misses some tank cannons & \textbf{Continued}, because how many real events we find (detection recall) is still incomplete \\
\addlinespace
Scan every few seconds for steady sounds + snap bangs via spectrum jumps (hop + flux onset) & Better vehicle/weather coverage, but cannons still often missed & \textbf{Continued} refining thresholds \\
\bottomrule
\end{tabularx}
\end{center}

\newpage
\subsection*{Making the vibration dense enough}
Reference hand-made map \texttt{videoplayback\_output\_haptic\_map[1].json}: about 71\% of short time windows vibrate, typical strength about 50/255, with multi-second continuous runs.

\begin{center}
\small
\begin{tabularx}{\textwidth}{@{}>{\raggedright\arraybackslash}p{4cm}>{\raggedright\arraybackslash}p{4cm}>{\raggedright\arraybackslash}X@{}}
\toprule
\textbf{What we tried} & \textbf{Outcome} & \textbf{Changed to / Continued?} \\
\midrule
Run A-D on the full clip as the quiet continuous layer (bed) without boosting quiet parts & Algorithm A often output silence on soft passages & \textbf{Changed to} boost quiet parts first (AGC level), then reshape loud/quiet from the source \\
\addlinespace
Mute bed whenever it was quiet relative to the loudest peak & Silenced most of algorithm D's rumble & \textbf{Changed to} mute only the quietest share of the track (percentile silence) \\
\addlinespace
Squeeze loud/quiet gap with a simple power curve (envelope-exponent compression) & Barely improved denseness for D & \textbf{Changed to} AGC level + loud/quiet curve taken from the source \\
\addlinespace
Average nearby gain values (mean-smooth compression) & Boosted bangs together with the quiet parts next to them & \textbf{Changed to} remember recent loud levels (peak-hold) while boosting \\
\addlinespace
Scale the bed by its single loudest sample (normalize bed by peak) & Quiet continuous layer too weak vs strong hits (accents) & \textbf{Changed to} scale bed by a high typical level, not the single peak \\
\bottomrule
\end{tabularx}
\end{center}

\section{Key takeaways}
\begin{itemize}[leftmargin=*,itemsep=2pt]
  \item Sound2Hap A-D stay as ports of the official scripts. We change context filtering, stitching, and level shaping instead of rewriting the algorithms.
  \item Finding event times alone is not enough, because vibration that only fires on detected events feels sparse next to a dense hand-made reference. Quiet continuous vibration (bed) plus stronger short hits (accents) is needed for thunder and rumble.
  \item How many real events we find (detection recall) on annotated tank footage is still incomplete (missed cannons). Better denseness does not replace better detection.
\end{itemize}

\end{document}
